% Standalone theorem/proof fragment. Uses amsmath, amssymb, amsthm.
\section{Polynomial-overhead robustification}
\label{sec:digitize}

\begin{proposition}
Suppose a width-$C$, depth-$L$ tanh network, with every weight bounded in
absolute value by $C$, implements $f:\{-1,1\}^d\to\{-1,1\}$ with strict
output margin $1/2$. Then $f$ has a tanh implementation of width and depth
polynomial in $(d,C,L)$ that tolerates simultaneous adversarial hidden-layer
activation perturbations of inverse-polynomial RMS size.
More precisely, there is a fan-in-two Boolean circuit for $f$ of size
$S=\operatorname{poly}(d,C,L)$. Put
\[
 W=\max\{8,d+S+1\},\qquad D=S+2.
\]
There is a width-$W$, depth-$D$ tanh network, all of whose weights have absolute
value at most $8$, such that every sequence of post-activation injections with
\[
 \|\zeta^s\|_2\le\frac18\quad(1\le s<D),\qquad
 |\zeta^{\rm out}|\le\frac1{8\sqrt W}
\]
preserves signed output margin strictly greater than $5/8$, on every Boolean
input. Its weighted norm is at most
\[
 \frac{64d(d+1)}W+192SW+W.
\]
At every width $N\ge W$, cloning preserves this guarantee at RMS radius
$1/(8\sqrt{2W})$ and increases weighted norm by at most a factor of two.
The construction allows polynomial depth and width overhead; it makes no
claim at the original depth $L$ or every original width $N\ge C$.
\end{proposition}

\begin{proof}
We first replace the analog computation by a finite Boolean computation.
Round every parameter to absolute precision $\eta=2^{-p}$. Evaluate each
hidden tanh to precision $\eta$ and clip its approximation into $[-1,1]$;
evaluate the readout to the same precision. If $E_s$ is the maximum coordinate
error after layer $s$, the first hidden layer satisfies
$E_1\le(d+2)\eta$, and the subsequent layers and readout satisfy
\[
 E_s\le C^2E_{s-1}+(C+1)\eta.
\]
This follows by writing each affine error as the original weights times
activation errors, plus weight errors times bounded approximate activations.
The clipping operation does not increase the approximation error.
It suffices to choose
\[
 2^{-p}\,[d+2+L(C+1)](C^2+1)^{L-1}\le\frac18.
\]
Thus $p=O(L\log(C+1)+\log(d+L+1))$ bits suffice, and the final sign is unchanged.
The rounded parameters are hardwired constants; no efficient description of
the original real numbers is required.

Here is an elementary polynomial-cost tanh evaluation algorithm. For a dyadic
argument $t$ with $|t|\ge p+2$, return its sign; the error is less than $2^{-p}$.
Otherwise set $x=2|t|\le2(p+2)$ and approximate $e^{-x}$ by its degree
$m=16(p+2)$ Taylor polynomial. Its error is at most
\[
 \frac{x^{m+1}}{(m+1)!}
 \le\left(\frac{ex}{m+1}\right)^{m+1}
 \le(e/8)^{m+1}<2^{-p-4}.
\]
The Taylor sum can be evaluated exactly using the common denominator
$2^{bm}m!$, where $b$ is the number of fractional bits of $x$.
All integers involved have polynomial bit length, even though some Taylor
terms are large and cancel. Round and clip the answer into $[0,1]$, then use
$(1-u)/(1+u)$ and oddness to obtain $\tanh t$.
This rational map has derivative of magnitude at most $2$ on $[0,1]$, so a
constant number of guard bits gives the required accuracy. Integer addition,
multiplication, comparison, and division have polynomial-size Boolean
circuits. Applying this algorithm throughout the rounded network gives a
Boolean circuit of size $S=\operatorname{poly}(d,C,L)$ for $f$.

We now simulate that circuit by noisy tanh gates. Express bits as $\pm1$, and
carry a constant $+1$ channel. Use
\[
 \begin{array}{ll}
 \mathrm{COPY}(a)=\tanh(8a),&\mathrm{NOT}(a)=\tanh(-8a),\\
 \mathrm{AND}(a,b)=\tanh(8(a+b-c)),&
 \mathrm{OR}(a,b)=\tanh(8(a+b+c)),
 \end{array}
\]
where $c$ is the constant channel, itself updated by $\tanh(8c)$.
Simplify gates with repeated or constant inputs beforehand, so each displayed
nonzero edge has magnitude $8$.
If all incoming signals are within $1/4$ of their ideal values, the signed
affine expression for AND and OR is at least $1-3/4=1/4$; COPY and NOT have
larger margin. Consequently each gate output is within
$1-\tanh2<1/8$ of its ideal value. Adding any coordinate perturbation of
magnitude at most $1/8$ leaves the error strictly below $1/4$.
This invariant holds at arbitrarily many layers, with no accumulated error.

Use one Boolean gate per hidden layer and copy all previously computed
signals forward. The first hidden layer initializes the input signals and
constant from $\tanh(8x_i)$, including $x_0=1$.
The $W$ channels accommodate all inputs, gate outputs, and the constant.
Unused channels have zero outgoing weights until used.
Each hidden-layer Euclidean noise budget $1/8$ bounds every coordinate
injection by $1/8$, so the invariant applies. Read out the desired final
signal with coefficient one. It has signed value greater than $3/4$; the
output injection is at most $1/8$, leaving signed margin greater than $5/8$.
The first layer has $d+1$ nonzero edges, each of the $S$ gate layers has at
most $3W$ nonzero edges, and the readout has one edge. This gives the stated
weighted norm.

Finally, at width $N\ge W$ use $k=\lfloor N/W\rfloor$ copies of each hidden
neuron, divide outgoing weights by $k$, and leave the remaining channels dead.
Clone-group means obey exactly the original computation with averaged
post-activation injections. Their Euclidean noise norm is at most
\[
 \frac1{8\sqrt{2W}}\sqrt{\frac Nk}\le\frac18,
\]
because $N/k\le2W$. Thus the same invariant and margin hold.
The first-layer norm does not increase, internal-layer norm is unchanged,
and the weighted readout norm increases by at most two.
\end{proof}

\paragraph{Interpretation.}
Inverse-polynomial internal robustness does not prevent representation of
any Boolean function already having a bounded-weight tanh implementation,
provided polynomial increases in depth and width are allowed. This is an
expressivity statement, not a claim that gradient descent produces the
constructed digital simulation. The point is that exponentially small analog
precision costs only polynomially many digital bits.
